\section{The universal relation group and interpolation}
\label{sec:depth}

Binary Goldbach moves determine both a monoid of monomial classes and
its group completion. We compute the latter on the odd primes and
deduce the interpolation formula of Theorem~\ref{main:prime-interpolation}.
We then prove that the finite relation groups stabilize, and use
localization to bound the depth and projective dimension of the two
truncations.

For the binary Goldbach ring $G=R/I$ on a finite support $\A$, let
$e_a$ denote the unit vector at $a$ in the free abelian group
$\mathbb Z^{\A}$. Set
\[
 L_{\A}=\langle e_a+e_b-e_{a'}-e_{b'}:a+b=a'+b',\ a,b,a',b'\in\A\rangle,
 \qquad U_{\A}=\mathbb Z^{\A}/L_{\A},
 \qquad r_{\A}=\operatorname{rank}_{\mathbb Z}U_{\A}.
\]
This is the universal ambient group for Freiman homomorphisms of
order two, as constructed by Konyagin--Lev \cite{KonyaginLev00}
and treated categorically by Blanco--Haghverdi \cite{BlancoHaghverdi26}.
Explicitly, every map $g:\A\to H$ to an abelian group that preserves
equal pair sums extends uniquely to a group homomorphism
$U_{\A}\to H$. We use this universal group to relate arithmetic
identities to localization and depth.

Write $[a]$ for the image of $e_a$. A monomial
$x_{\boldsymbol a}=x_{a_1}\cdots x_{a_d}$ determines the element
$[a_1]+\cdots+[a_d]$ of $U_{\A}$. The classes of monomials modulo binary moves form a commutative
monoid under multiplication. Its monoid algebra is $G$, by
Proposition~\ref{setup:components}. The group $U_{\A}$ is
its group completion, in which the class of each variable has an inverse.

\subsection{The integral group of prime relations}

For an infinite support $\A$, the same construction uses finitely
supported integer exponent vectors. Write $U_\infty=U_{\Podd}$ for the
group on all odd primes. Its relations are
$[p]+[q]=[p']+[q']$ whenever $p+q=p'+q'$. To compute
this group, first omit $3$. In a pair of primes greater
than three, the number $j$ of entries congruent to one modulo three
is determined by the sum, since the three possibilities $j=0,1,2$ give
distinct values of $2j-2\bmod3$. For example, $5+19=7+17$ replaces one prime in each residue class
by another in the same class. Every binary move among primes greater
than three preserves both counts. To prove the converse in the group, we construct
pair-sum identities among primes and cancel the auxiliary generators.

\begin{lemma}\label{depth:prime-pair-conversion}
Let $\boldsymbol p=(p_1,\ldots,p_d)$ and $\boldsymbol q=(q_1,\ldots,q_d)$
be lists of primes in $\Podd_{>3}=\Podd\setminus\{3\}$ with the same sum
and the same counts in each residue class modulo three. There is a
(possibly empty) list $\boldsymbol r$ of primes in $\Podd_{>3}$ such
that binary moves on this support connect $x_{\boldsymbol r}x_{\boldsymbol p}$
to $x_{\boldsymbol r}x_{\boldsymbol q}$. In particular,
$\sum_i[p_i]=\sum_i[q_i]$ in $U_{\Podd_{>3}}$.
\end{lemma}
\begin{proof}
For $d\le1$ the monomials coincide. For $d=2$ the two monomials
either coincide or are joined by a single binary move.
Assume $d\ge3$, and reorder the lists so that
$p_i\equiv q_i\pmod3$ for each $i$.
Put $\Delta_0=0$ and $\Delta_i=\sum_{j=1}^i(p_j-q_j)$.
Each $\Delta_i$ is divisible by six, and the equality of the sums gives
$\Delta_d=0$.

Choose integer parameters $a_1,\ldots,a_d,c$. With indices read cyclically
on $a_i$ and $b_i$, define
\[
 b_i=a_{i+1}+c+\Delta_i,\qquad r_i=a_i+b_i-p_i .
\]
For every choice of the parameters we have the integer identities
\[
 p_i+r_i=a_i+b_i,\qquad q_i+r_i=a_{i+1}+b_{i-1}.
\]
The second follows from $\Delta_i-\Delta_{i-1}=p_i-q_i$.
We seek one parameter tuple at which the $3d$ values $a_i,b_i,r_i$
are all primes greater than three. At such a tuple the auxiliary
integer $c=b_i-a_{i+1}-\Delta_i$ is even.

We obtain simultaneous primes from the theorem on linear equations
in primes of Tao and Ter\"av\"ainen \cite[Theorem~1.6]{TT}.
Only its qualitative asymptotic is needed here; this was already
available from the work of Green--Tao and Green--Tao--Ziegler
\cite{GT10,GT12,GTZ12}, as recalled immediately after that theorem.
The linear parts $e_i$, $e_{i+1}+e_c$, $e_i+e_{i+1}+e_c$ are pairwise
nonproportional, since $d\ge3$. Modulo $2$ take $c=0$, $a_i=1$.
All the forms are odd. Modulo $3$ take $c=1$, $a_i=1$. Then $b_i=2$ and
$r_i=-p_i\ne0\bmod3$. At a prime $s\ge5$ set $c=0$ modulo $s$ and choose
the $a_i$ successively. First choose $a_1\ne0$; since $\Delta_d=0$,
this also makes $b_d=a_1\ne0$. For $2\le i\le d$, choose $a_i$ avoiding
\[
 0,\qquad -\Delta_{i-1},\qquad
 p_{i-1}-a_{i-1}-\Delta_{i-1},
\]
so that $a_i,b_{i-1},r_{i-1}$ are nonzero.
At the final choice also avoid $p_d-a_1$, ensuring
$r_d=a_d+a_1-p_d\ne0$. Each choice excludes at most four residues,
so it is possible even at $s=5$. Thus all local factors are positive.
For large $s$, counting pairwise intersections of the corresponding
hyperplanes gives normalized local factors $1+O_d(s^{-2})$, so
their product is positive.

Take all $a_i,c$ in $[Y,2Y]$, with the two lists fixed. Every form
is positive and comparable to $Y$ for large $Y$. This convex box lies
in $[-2Y,2Y]^{d+1}$, so we apply the theorem with $3d$ forms in $d+1$
variables and scale $N=2Y$. Its affine constants are fixed with the
two lists and hence satisfy the required $O(N)$ bound. The theorem gives
a von Mangoldt weighted count $\gg_{\boldsymbol p,\boldsymbol q}Y^{d+1}$. If one form is a
proper prime power, its value has $O(\sqrt Y\log Y)$ possibilities
and the other $d$ free variables have $O(Y^d)$ choices. Including all
logarithmic weights gives $o(Y^{d+1})$ for these points. Hence some
integer tuple gives prime values greater than three for all $3d$ forms.
For this tuple put $\boldsymbol a=(a_1,\ldots,a_d)$,
$\boldsymbol b=(b_1,\ldots,b_d)$ and $\boldsymbol r=(r_1,\ldots,r_d)$.
Replacing each pair $(p_i,r_i)$ by $(a_i,b_i)$ gives a path from
$x_{\boldsymbol r}x_{\boldsymbol p}$ to $x_{\boldsymbol a}x_{\boldsymbol b}$.
Replacing $(q_i,r_i)$ by $(a_{i+1},b_{i-1})$ gives a path from
$x_{\boldsymbol r}x_{\boldsymbol q}$ to the same monomial, by cyclic indexing.
This proves the assertion about moves. Passing to the group allows us to
cancel $\sum_i[r_i]$.
\end{proof}

\begin{theorem}[The integral group of binary prime relations]
\label{depth:prime-pair-group}
The group on the primes greater than three is free of rank three.
On all odd primes the group $U_\infty$ is free of rank two,
with
\[
 [p]=[5]+\frac{p-5}{2}\bigl([7]-[5]\bigr)
 \qquad(p\in\Podd).
\]
\end{theorem}
\begin{proof}
Write $A=[5]$, $B=[7]$, and $D=[11]-[5]$. Every prime $p>3$ has
the form $s+6j$ with $s=5$ or $7$ and $j\ge0$.
Apply Lemma~\ref{depth:prime-pair-conversion} to $p$ together with $j$
copies of $5$, and $s$ together with $j$ copies of $11$. It gives
$[p]=[s]+jD$. For example, $5+17=11+11$ gives $[17]=[5]+2D$.
Thus $A,B,D$ generate the group on primes greater
than three. The residue counts distinguish $A$ and $B$. Both counts
vanish on $D$, which has weight six. These three elements are
therefore integrally independent.

Now adjoin $3$. The relations $3+7=5+5$ and $3+11=7+7$ give
$[3]=2A-B$ and $D=3(B-A)$. Thus the step between primes in the
same residue class is $3(B-A)$ per six units of weight, and
$[7]-[5]=B-A$ accounts for the two-unit step between $5$ and $7$.
Substituting in $[p]=[s]+jD$ gives the asserted formula for every
$p>3$, and $[3]=2A-B$ gives it also at $3$. The generators $A$ and
$B-A$ have independent vectors of Goldbach and monomial degrees
$(5,1)$ and $(2,0)$, so
they freely generate.
\end{proof}

To deduce Theorem~\ref{main:prime-interpolation}, observe that any map
$g:\Podd\to H$ preserving equal pair sums extends to a homomorphism
$U_\infty\to H$ by $[p]\mapsto g(p)$. Since
$[7]-[5]=[5]-[3]$, the formula of Theorem~\ref{depth:prime-pair-group}
becomes $[p]=[3]+((p-3)/2)([5]-[3])$. Applying the homomorphism gives
the interpolation formula in $H$.

For complex-valued maps the formula is equivalent to $g(p)=Ap+B$
on the odd primes. If $f(p+q)=g(p)+g(q)$ for odd primes $p,q$, it
also gives $f(n)=An+2B$ on $\Podd+\Podd$. This proves the odd-prime
form of the binary affine-classification conjecture in
\cite[Conj.~1]{KP15}. The corresponding theorem for sums of at least
three primes is proved there; Park \cite[p.~6, after Thm.~3.7]{Park25}
still presents the binary analogue as conjectural. Our integral
computation additionally allows arbitrary abelian targets, including
groups with torsion.

The odd-prime restriction must be distinguished from the all-primes
wording of \cite[Conj.~1]{KP15}. Every equal-pair-sum relation involving
$2$ is trivial: an odd sum $2+p$ has only the unordered prime
representation $\{2,p\}$, and $4$ has only $2+2$.
Thus, if $U_{\{2\}\cup\Podd}$ denotes the universal group on all primes,
\[
 U_{\{2\}\cup\Podd}=\mathbb Z[2]\oplus U_\infty
                    \cong\mathbb Z^3.
\]
The value $g(2)$ is independent of the values on the odd primes, so
pair-sum preservation alone cannot imply one affine formula on all
primes. For example, $g(2)=1$ and $g(p)=0$ for every odd prime $p$
preserves equal pair sums but is not affine on the full prime set.
The comparison with \cite[Conj.~1]{KP15} is therefore with its
odd-prime form. The shifted conjecture addressed next already asks
for an odd-prime conclusion.

\begin{corollary}[K\'atai--Phong interpolation]\label{depth:KP-interpolation}
Fix positive integers $a,b$. Let $f,g$ be functions on the positive
integers with values in an abelian group, and suppose that
\[
 f(p+q+a+b)=g(p+a)+g(q+b)
\]
for all odd primes $p,q$. Then for every odd prime $p$,
\[
 g(p+a)=\frac{p-3}{2}\,g(5+a)-\frac{p-5}{2}\,g(3+a).
\]
In particular this proves Conjecture~6 of K\'atai and Phong
\cite[p.~235]{KP15Problems}, whose hypothesis imposes the displayed
identity for all primes and whose conclusion concerns the odd primes.
\end{corollary}
\begin{proof}
Put $g_a(p)=g(p+a)$ and $g_b(p)=g(p+b)$. Interchanging $p$ and $q$ gives
$g_a(p)-g_b(p)=g_a(q)-g_b(q)$. Thus $g_a(p)+g_a(q)$ differs from $f(p+q+a+b)$
by a constant and depends only on $p+q$. Apply
Theorem~\ref{main:prime-interpolation} to $g_a$ and collect the
coefficients of $g_a(3)$ and $g_a(5)$.
\end{proof}

\subsection{Finite stabilization}

Put $\A_m=\{p\le m:p\text{ odd prime}\}$ and $U_m=U_{\A_m}$.

The existence of a relation involving only smaller primes was stated by
K\'atai--Phong \cite[Theorem~3]{KP15}: every sufficiently large prime
is a sum of two smaller primes minus a third smaller prime. We use this
arithmetic fact to obtain integral stabilization below. For completeness,
we give a short proof with balanced ranges from the classical almost-all
Goldbach estimate. The balanced form will also distinguish the two
truncations.

Goldbach questions also enter the numerical semigroups generated by
primes above a threshold: Hellus--Rechenauer--Waldi
\cite{HellusRechenauerWaldi20} study their Frobenius numbers and
minimal generators. The relations used here preserve both the sum
and the number of summands.

Write $Y$ for the interval scale and $n$ for Goldbach degree. Set
\[
 \begin{aligned}
 T_{2,Y}(n)&=\#\{(p,p'):p,p'\text{ odd primes},\ p+p'=n,\ Y<p,p'\le2Y\},\\
 R_{2,Y}(n)&=\sum_{\substack{u+v=n\\Y<u,v\le2Y}}\Lambda(u)\Lambda(v),
 \end{aligned}
\]
where $u,v$ range over integers and $\Lambda$ is the von Mangoldt function.
Let $V_{2,Y}(n)$ count all integer pairs in the same range, without weights,
and let $\mathfrak S(n)$ denote the Goldbach singular series.
The analytic estimate recalled in the introduction of
Bazzanella--Languasco \cite[eq.~(2)]{BL00}, with their parameters
$H=1$ and $N=Y$, gives the single-target estimate
$R_{2,Y}(n)\sim V_{2,Y}(n)\mathfrak S(n)$ for all but
$O_C(Y/(\log Y)^C)$ even integers $n\in(5Y/2,7Y/2]$, for every fixed $C>0$.
In this interval the first factor is $\gg Y$
and the second is bounded below by a positive constant. The contribution
involving a proper prime power is $O(Y^{1/2}(\log Y)^3)=o(Y)$: there are
$O(Y^{1/2}\log Y)$ such prime powers up to $2Y$, and each weight is
at most $(\log(2Y))^2$. Thus $T_{2,Y}(n)>0$ for every nonexceptional
even target and sufficiently large $Y$, giving a representation by primes
\[
 n=p+p',\qquad Y\le p,p'\le2Y.
\]
In fact the primes can be chosen strictly greater than $Y$. The parameter
$m$ remains the prime cutoff: $2Y\le m$ ensures that both primes belong to
$\A_m=\{p\le m:p\text{ odd prime}\}$ for large $Y$. We use the recalled
almost-all estimate, not the main short-interval theorem of \cite{BL00}.

\begin{lemma}[Balanced lower-prime relations]\label{depth:moving-anchor}
There is an integer $q_0$ such that every prime $q>q_0$ admits odd primes
$p,p',p''<q$ such that
\[
 p+q=p'+p'',\qquad q/4<p<q/2,\qquad (q-1)/2<p',p''\le q-1.
\]
\end{lemma}
\begin{proof}
Put $Y=(q-1)/2$. By the prime number theorem, the interval $(q/4,q/2)$
contains asymptotically $q/(4\log q)$ prime anchors $p$. The distinct
even targets $n=p+q$ lie in $(5Y/2,7Y/2]$ for large $q$. With a fixed
$C>1$, the exceptional set in the preceding estimate has
$O_C(q/(\log q)^C)=o(q/\log q)$ elements. At least one target
therefore has a genuine prime representation $p'+p''$ with
$Y<p',p''\le2Y=q-1$, as required.
\end{proof}


Return to the odd-prime supports $\A_m$.
Inclusion of supports induces maps
$U_m\to U_{m'}$, where
$U_m=U_{\A_m}$ and $m\le m'$.

\begin{theorem}\label{depth:integral-stabilization}
Let $U_\infty$ be the group on all odd primes modulo the
equal-pair-sum relations. For every sufficiently large $m$, the natural
map $U_m\to U_\infty$ is an isomorphism.
\end{theorem}
\begin{proof}
Choose $m_0>q_0$, so the prefix $\A_{m_0}$ contains all odd primes
at most $q_0$. For each
subsequent prime $q$, Lemma~\ref{depth:moving-anchor} gives
\[
 [q]=[p']+[p'']-[p]
\]
using only earlier labels. Hence every subsequent transition is
surjective, and all later groups are quotients of one fixed finitely
generated abelian group $U_{m_0}$. The kernels of the maps from $U_{m_0}$ form an
ascending chain of subgroups. It stabilizes because $U_{m_0}$ is a Noetherian
$\mathbb Z$-module. Stabilization of the kernels proves stabilization
of the groups and their transition maps, including their torsion.
The direct limit is $U_\infty$, since every defining relation
uses finitely many prime labels.
\end{proof}

\subsection{Depth and projective dimension}

For the binary odd-prime ring, use the cutoff notation of
Section~\ref{setup:resolution-construction}, with
$\A_m=\{p\le m:p\text{ odd prime}\}$ and $R_m=R^{<m}$.
Depth is taken at the homogeneous maximal ideal unless another
localization is specified. By \eqref{setup:AB-cutoffs}, depth measures
the difference between $|\A_m|$ and projective dimension. We bound
this difference for both truncations.

Both bounds use the same homological observation: localization preserves
a free resolution, although it need not preserve minimality. Thus the
projective dimension at the all-ones point is a lower bound for the
projective dimension over $R$. A dimension bound there, together with
Auslander--Buchsbaum, will give the desired bound.

\begin{lemma}[Triangular elimination]\label{depth:triangular}
Let $\kk$ be a field, let $\A$ be a finite ordered support with least label $a_0$, and put
$R=\kk[x_a:a\in\A]$. Suppose
$\mathfrak a\subseteq\mathfrak q=\langle x_a-1\mid a\in\A\rangle$.
Let $\mathcal B\subseteq\A\setminus\{a_0\}$, and suppose that for each
$b\in\mathcal B$ there are $a',a''\in\A$, depending on $b$, such that
\[
 x_{a_0}x_b-x_{a'}x_{a''}\in\mathfrak a,\qquad a',a''<b.
\]
Then $\pd_R(R/\mathfrak a)\ge|\mathcal B|$. If $\mathfrak a$ is homogeneous in monomial degree,
then $\operatorname{depth}(R/\mathfrak a)\le|\A|-|\mathcal B|$.
\end{lemma}
\begin{proof}
Put $M=R/\mathfrak a$, and let $\mathfrak b$ be generated by the displayed
relations. For recursive substitution, only $x_{a_0}$ needs to be inverted:
each equation gives $x_b=x_{a_0}^{-1}x_{a'}x_{a''}$. In increasing order
of $b\in\mathcal B$, we therefore solve for $x_b$ in terms of earlier
variables, obtaining
\[
 R[x_{a_0}^{-1}]/\mathfrak b\cong
 \kk[x_{a_0},x_{a_0}^{-1},x_c:c\in\A\setminus(\mathcal B\cup\{a_0\})],
\]
whose dimension is $|\A|-|\mathcal B|$.

We now localize further at $\mathfrak q$ to apply Auslander--Buchsbaum
over the regular local ring $R_{\mathfrak q}$, whose depth and dimension
are $|\A|$. Every $x_a$ becomes a unit there, since $x_a\notin\mathfrak q$,
but these additional inverses were not needed for substitution.
The inclusions $\mathfrak b\subseteq\mathfrak a\subseteq\mathfrak q$
ensure that $M_{\mathfrak q}\ne0$ and
$\dim M_{\mathfrak q}\le|\A|-|\mathcal B|$.
Localization preserves a free resolution, so Auslander--Buchsbaum gives
\[
 \pd_R M\ge\pd_{R_{\mathfrak q}}M_{\mathfrak q}
 =|\A|-\operatorname{depth}M_{\mathfrak q}\ge|\mathcal B|.
\]
If $\mathfrak a$ is homogeneous, the graded minimal resolution of $M$
stays minimal at the homogeneous maximal ideal
$\mathfrak m=\langle x_a\mid a\in\A\rangle$. Since $R_{\mathfrak m}$
also has depth $|\A|$, applying Auslander--Buchsbaum there gives the
depth comparison and bound explicitly:
\[
 \begin{aligned}
 \operatorname{depth}M
 &=|\A|-\pd_R M\\
 &\le|\A|-\pd_{R_{\mathfrak q}}M_{\mathfrak q}
   =\operatorname{depth}M_{\mathfrak q}\\
 &\le\dim M_{\mathfrak q}\le|\A|-|\mathcal B|.
 \end{aligned}\qedhere
\]
\end{proof}

For a Goldbach relation the label identity is $a_0+b=a'+a''$;
all three labels $a_0,a',a''$ are smaller than $b$.

\begin{proposition}\label{depth:relation-group-bound}
For every finite nonempty support $\A$ and every field $\kk$,
\[
 \operatorname{depth}G\le r_{\A},\qquad
 \pd_{R}G\ge|\A|-r_{\A}.
\]
\end{proposition}
\begin{proof}
Put $w=x_{\A}=\prod_{a\in\A}x_a$. Inverting $w$ makes each $x_a$
invertible. In the Laurent polynomial ring write
$x^\alpha=\prod_{a\in\A}x_a^{\alpha(a)}$ for $\alpha\in\mathbb Z^{\A}$,
allowing negative exponents. Each binary relation becomes
$x^{e_a+e_b-e_{a'}-e_{b'}}=1$. Since $x^{\alpha+\beta}=x^\alpha x^\beta$
and $x^{-\alpha}=(x^\alpha)^{-1}$, these relations identify exactly
the Laurent monomials whose exponents lie in the same coset of $L_{\A}$.
Thus
\[
 G[w^{-1}]\cong\kk[U_{\A}],
\]
with one basis element for each coset and multiplication given by
addition in $U_{\A}$. This is the Laurent form of the lattice
ideal construction \cite[Theorem~2.1]{ES96}, \cite[\S7.1]{MS05}.
By the structure theorem for finitely generated abelian groups,
$U_{\A}\cong\mathbb Z^{r_{\A}}\oplus T_{\A}$, with $T_{\A}$ finite.
Its group algebra is a Laurent polynomial algebra over the nonzero
finite-dimensional algebra $\kk[T_{\A}]$, and hence has dimension $r_{\A}$.
This remains true when the characteristic divides $|T_{\A}|$.

All defining binomials vanish at the all-ones point, so the quotient
localized at $\mathfrak q=\langle x_a-1\mid a\in\A\rangle$ is nonzero and has dimension
at most $r_{\A}$. Repeating the final argument of
Lemma~\ref{depth:triangular}, local
Auslander--Buchsbaum gives $\pd_{R}G\ge|\A|-r_{\A}$, and graded
Auslander--Buchsbaum gives the asserted depth bound.
\end{proof}

\begin{corollary}[Depth and total Betti numbers]
\label{depth:bounded-depth}
For all sufficiently large $m$, over every field,
\[
 \delta^{<m}\le2,\qquad
 |\A_m|-2\le\rho^{<m}\le |\A_m|.
\]
Moreover the total Betti numbers satisfy
$\beta^{<m}_{j,*,*}\ge\binom{|\A_m|-2}{j}$ for $0\le j\le |\A_m|-2$.
\end{corollary}
\begin{proof}
By Theorems~\ref{depth:integral-stabilization} and
\ref{depth:prime-pair-group}, $U_m\cong\mathbb Z^2$ for all
sufficiently large $m$. Hence $L_{\A_m}$ is a direct summand of
$\mathbb Z^{|\A_m|}$ of rank $|\A_m|-2$. After inverting all variables,
extend a basis of $L_{\A_m}$ to a basis of the exponent lattice.
In the resulting Laurent coordinates the ideal is generated by
$y_1-1,\ldots,y_{|\A_m|-2}-1$, a regular sequence. At the point with
all coordinates one its Koszul resolution is minimal. The graded
minimal resolution over $R_m$ localizes to a free resolution there,
and tensoring with the residue field gives a complex whose $j$th term
has dimension $\beta^{<m}_{j,*,*}$ and whose $j$th homology is the
local Tor group of dimension $\binom{|\A_m|-2}{j}$. Since homology
cannot have larger dimension than its chain group, this gives the
termwise Betti bounds. The
projective dimension bounds now follow from the last nonzero local
rank and Hilbert's syzygy theorem. Graded Auslander--Buchsbaum gives
the depth bound.
\end{proof}

In particular $\rho^{<m}=|\A_m|-O(1)$, with
$|\A_m|=\pi(m)-1$ for integers $m\ge3$. The auxiliary primes in
Lemma~\ref{depth:prime-pair-conversion} may exceed the given cutoff,
so the isomorphism in Theorem~\ref{depth:integral-stabilization} is
essential to this conclusion. Its proof leaves the numerical onset of
stabilization open. Corollary~\ref{depth:verified-range} below verifies
the same rank-two group and depth bound throughout an explicit finite
range, without identifying the onset in the eventual theorem.

\subsection{The relation-weight cutoff}

The relation $p+q=p'+p''$ in Lemma~\ref{depth:moving-anchor} can have weight
larger than $m$ when $q$ is near the cutoff. It therefore does not give
the same uniform depth bound for $G^{\le m}$. A fixed anchor yields a
weaker but still asymptotically sharp projective-dimension estimate.

For an integer $m\ge6$, set
\[
 \begin{aligned}
 E_m&=\{p\in\A_m:p+3\ne q+r\text{ for all primes }q,r>3\},
       &e_m&=|E_m|,\\
 u_m&=\#\{p\in\A_m:p>m-3\}.&&
 \end{aligned}
\]
Repeated summands are allowed, and $u_m\le2$ since $(m-3,m]$ contains at most two odd integers.

\begin{proposition}\label{depth:fixed-anchor}
Over every field,
\[
 \delta^{<m}\le e_m,
 \qquad\delta^{\le m}\le e_m+u_m.
\]
\end{proposition}
\begin{proof}
For $p\notin E_m$, choose $p+3=q+r$ with $q,r>3$ prime. Then
$q,r\le p-2<p$, so $x_3x_p-x_qx_r$ is a triangular relation in the
variable cutoff. It is present in the relation cutoff when $p\le m-3$.
All binomials vanish at the all-ones point. Apply
Lemma~\ref{depth:triangular}, omitting at most $u_m$ relations in the
second case.
\end{proof}

\begin{corollary}[A verified range of variable cutoffs]
\label{depth:verified-range}
Put $B=20\,000\,000$. Over every field, for every integer $3\le m\le B$,
\[
 \depth G^{<m}=\delta^{<m}\le2.
\]
For $5\le m\le B$, moreover, $U_m\cong\mathbb Z^2$ and
\[
 \beta^{<m}_{i,*,*}\ge\binom{|\A_m|-2}{i}
 \qquad(0\le i\le|\A_m|-2).
\]
\end{corollary}
\begin{proof}
The exhaustive sieve in
\path{supplement/computations/prime-pair-depth-check.py} verifies that
each of the $1\,270\,604$ primes $7\le p\le B$ has
$p+3=q+r$ with primes $q,r>3$. Thus $E_m=\{3,5\}$ for
$6\le m\le B$, and Proposition~\ref{depth:fixed-anchor} gives the
depth bound. The cases $m=3,4,5$ are polynomial rings in at most two
variables.

Every verified pair has $q,r\le p-2$, so it lies in the same prime
prefix as $p$. Induction on $p$ in the relation group, using
$[p]=[q]+[r]-[3]$, gives
\[
 [p]=(1-j)[3]+j[5],\qquad j=(p-3)/2.
\]
Hence $[3],[5]$ generate $U_m$. The map
$[p]\mapsto(1-j,j)\in\mathbb Z^2$ respects every equal-pair-sum
relation and sends these two generators to the standard basis.
They are therefore independent, proving $U_m\cong\mathbb Z^2$.
The localization argument in the proof of
Corollary~\ref{depth:bounded-depth} now gives the total Betti bounds
over every field.
\end{proof}

This computation does not determine whether its verified range
overlaps the range of the eventual theorem, and it does not prove
the depth bound at every cutoff. The fixed-anchor condition is
sufficient for the bound; no necessity is asserted. For the
relation-weight cutoff the same verification gives only
$\delta^{\le m}\le2+u_m\le4$ for $6\le m\le B$, since a relation
with anchor $3$ and label $p$ has weight $p+3$.

\begin{theorem}\label{depth:relation-cutoff-pd}
For every fixed $C>0$, over every field,
\[
 0\le\delta^{\le m}
   =|\A_m|-\rho^{\le m}
   \ll_C m(\log m)^{-C}.
\]
Consequently $\rho^{\le m}/|\A_m|\to1$.
\end{theorem}
\begin{proof}
Let $E^*(X)$ count the even integers at most $X$ having no representation
as a sum of two primes greater than three. The balanced almost-all
estimate above implies
\[
 E^*(X)\ll_C X(\log X)^{-C}\qquad(C>0).
\]
To see this with the required restriction on summands, cover
$[\sqrt X,X]$ by the intervals $(5Y/2,7Y/2]$ at geometric scales
with ratio $4/3<7/5$. Their lengths sum to $O(X)$ and their logarithms
are comparable to $\log X$. At each scale, nonexceptional targets
have both prime summands greater than three. Integer rounding does not
destroy the overlap for large scales. The remaining initial interval
contributes $O(\sqrt X)$, which is absorbed in the stated bound.

The injective map $p\mapsto p+3$ sends $E_m$ into the exceptional set
counted by $E^*(m+3)$. Thus $e_m\ll_C m(\log m)^{-C}$. Apply
Proposition~\ref{depth:fixed-anchor} and $u_m\le2$. Finally,
$|\A_m|\sim m/\log m$, so taking $C>1$ proves the limiting ratio.
\end{proof}

The theorem does not show bounded depth for the relation cutoff or
characteristic-independent exact depth at each finite cutoff. For the
variable cutoff, Corollary~\ref{depth:bounded-depth} is stronger than
this limiting ratio.


The group computation also describes the comparison with the toric
quotient after localization. On the full odd-prime support, two
monomials of the same bidegree become equal after multiplication by
some common monomial. Appendix~\ref{app:saturation} proves the
squarefree-multiplier refinement, studies saturation by a single
prime variable, and gives examples where cancellation fails. It also
develops explicit dual Koszul cocycles and the stronger infinite
conjecture associated with them.
